% Ad Familiares 3.1 --- headnote
% ad-familiares-03-01, late February 53 BC, written at Rome

\textit{Cicero to Appius Claudius Pulcher, proconsul of Cilicia,
written at Rome late in the winter of 53 BC and placed first
in book 3 of the \textit{Familiares}, which collects the entire
correspondence between the two men. Appius is the elder brother
of Publius Clodius, the tribune who had driven Cicero into
exile in 58 and continues to be the most dangerous of his
enemies in the city; he is also the man whom Cicero is destined
to succeed in Cilicia in 51, when the lex Pompeia of 52 will
press the unwilling consular into a provincial command he has
spent his career avoiding. The letter accordingly opens an
exchange whose tone is among the most studied in the corpus:
formal, polished, full of the courtesies that allow two men
linked by deep mutual reserve to do necessary business without
acknowledging what stands between them. The vehicle is a
freedman, Phania, who is carrying the real news of affairs at
Rome to Cilicia and to whom Cicero defers for the substance; he
himself takes up only the rituals of goodwill.}

\textit{The set-piece is the Minerva joke at the close of \S 1.
Appius, like his line, had a familial cult of Minerva; Cicero,
playing along, promises that if the Minerva by whose favour he
means to repay Appius's good offices turns out to be borrowed
from Appius's own household gods, he will call her not only
\textit{Polias} --- the standard Athenian epithet, ``of the
city'' --- but \textit{Appias}, a coinage in Appius's honour.
The whole letter is at this pitch. \S 2 dispatches the second
freedman, Cilix, with the practised graciousness of a man
turning a recommendation into the appearance of friendship;
\S 3 commends the jurisconsult L. Valerius in a half-teasing
formula. The Perseus dateline (\textit{Romae vel ex. a. 701
\textnumero{} vel in. 702}) places the letter on the cusp
between late 53 and early 52; Shackleton Bailey settles on
late February 53 BC. Two years later Cicero will be writing
from Cilicia about the wreckage Appius left behind him, but
here the tone is still that of the careful opening move.}
